\section{Divide and Conquer}

\begin{definition}

\textbf{Divide and Conquer } refers to the strategy of splitting a problem into smaller subproblems,
                solving them recursively, and combining the results.
            
\end{definition}

Problems of DC with manual fork/join:
\begin{itemize}
    \item Poor reusability of code
    \item Commit to a number of threads a priori and can not dynamically adjust.
    \item Subproblems are not automatically assigned to a thread.
\end{itemize}


\begin{remark}
Some manual fork/join fixes are: sequential-cutoffs or running one
subproblem on the thread itself, instad of creating two new threads.
\end{remark}

\subsection{Executor Service}

The Executor Service framework provides the way to schedule and run tasks on a thread-pool automatically.

\begin{definition}
\textbf{Runnable}: Task without return value. (Passed to \texttt{execute()} or \texttt{submit()}).
\textbf{Callable}: Task returning \texttt{Future<T>}. (Passed to \texttt{submit()}).
\end{definition}

\begin{remark}
\textbf{Pool Types}: \texttt{newFixedThreadPool(n)} (keeps $n$ threads), \texttt{newCachedThreadPool()} (scales dynamically).
\end{remark}

Pros and Cons of Executor service:
\begin{itemize}
    \item Not suited for: Recursive subproblems that depend on each other and have a deep fork-tree. (run out of threads)
    \item Suited for: flat structures that can run independently in parallel.
\end{itemize}

\subsection{Fork-Join}

Manages a pool of threads to solve recursive tasks using a work-stealing strategy, where threads have local task-queues
which they work on and if empty steal from other threads.

\begin{definition}
\textbf{RecursiveTask}: Returns a result (like Callable).
\textbf{RecursiveAction}: Does not return a result (like Runnable).
\textbf{Key Methods}: \texttt{fork()} (async push), \texttt{join()} (wait, but steals work while waiting), \texttt{compute()} (main logic).
\end{definition}

\begin{codeexample}
class MyTask extends RecursiveTask<Integer> {
    int start, end;
    MyTask(int start, int end) { this.start = start; this.end = end; }
    
    protected Integer compute() {
        if (end - start <= SEQUENTIAL_CUTOFF) { 
            return baseCase(); 
        } else {
            int mid = (start + end) / 2;
            MyTask left = new MyTask(start, mid);
            MyTask right = new MyTask(mid, end);
            
            left.fork(); // Async execute left half
            int rightRes = right.compute(); // Optimize: compute right half now
            int leftRes = left.join();      // Wait for left half
            
            return leftRes + rightRes;
        } 
    }
}
\end{codeexample}

\begin{important}

Make sure that every task has enough work to do (aprox. 100-10'000 ops.) (else the overhead will be larger than the parallel speedup gain)
            
\end{important}

\subsection{Task graphs}

The execution of parallel programs using Fork-Join can be represented as a Task-Graph (DAG).

\begin{definition}
\textbf{Span} $T_{\infty}$: execution time with unlimited hardware (Sum of longest dependent path in DAG).
Max speedup is limited by $T_1 / T_{\infty}$ (Amdahl's Law).
\end{definition}

        